\chapter{The Hard Cases}
\label{sec:hard}

Each case below is stated, run through the rules, and marked \emph{solved}, \emph{solved with a
cost}, or \emph{open}.

\section{Closures}
\label{sec:hard-closures}

\paragraph{A closure that reads a captured list, called many times.}
\begin{Verbatim}
total = List.foldl xs 0 (λx acc → acc + List.length ys)    -- reads ys, called n times
zs = List.set ys 0 7                                       -- then edits ys
\end{Verbatim}
The lambda holds \code{ys} under $\env$ and only reads it, so it is not once. \code{foldl} calls it many
times and does not keep it, and the edge check passes. The closure holds \code{ys} while it is live,
which is during the \code{foldl} call only, so at \code{List.set ys 0 7} the holder is dead.
\emph{Solved.}

The reason it is solved is that beni's lambdas are in front of the checker. In type-based usage
analyses the capture counts as a use of the capture's \emph{type}: in Wansbrough and Peyton Jones's
usage analysis, ``all free variables of an abstraction have at least the usage of the abstraction
itself'' \citep[\S6.1]{wansbrough1999once}. Effect- and capture-based systems close the case with
declarations on the API side---Deng et al.'s use and kill effects \citep{deng2025capture}, and
Capybara's markers on mutating and consuming methods \citep[\S6.1]{xu2026capybara}, although Capybara
infers read-only permissions from capture sets---and Affe with an explicit shared borrow \code{\&x}
\citep[\S3.3]{radanne2020kindly}.

\paragraph{A closure that consumes a capture, called many times.}
\begin{Verbatim}
ys2 = List.map xs (λx → List.push acc x)                    -- consumes a capture, called n times
\end{Verbatim}
The lambda writes its captured \code{acc}, so it is once: a cell, whose every call is a demand. \code{map}
calls it many times; the edge ``a once lambda requires $\mathit{calls} \le 1$'' fails, and the program is
\emph{rejected}, correctly: the first call writes \code{acc} and returns it, the second writes the same
array, and the first element of the result would change under the second call. The message names the
lambda, the capture and \code{map}'s call count, and offers a fold (\S\ref{sec:errors-examples}). This
is OxCaml's lock rule for closures, under which a closure that consumes a captured value must be
\code{once} \citep[\S3.3]{lorenzen2024oxidizing}, and Futhark's refusal to let the function of a
\code{map} consume an array bound outside it, since ``that would imply the array is consumed once for
every iteration'' \citep[\S3.3]{henriksen2017futhark}. A once closure that \emph{is} called at most
once---\code{Task.scope (λscope → ⋯ List.push acc x ⋯)}, \code{Result.andThen r (λx → List.push acc x)}---passes.
\emph{Solved}, with a cost: a once closure passed to a $\mathit{many}$ position is an error even when the
program would in fact call it once for reasons the checker cannot see, for example a \code{List.map} over
a one-element list.

\paragraph{A closure that leaves its binding.} An earlier version of these rules attached capture and
once-ness to the binding that named a lambda, and lost both when the closure was returned, stored or
captured by another closure: a function returning \code{λ⊤ → List.length xs} made its caller's
\code{xs} look unheld, and a function returning \code{λx → List.push acc x} could be called twice through
its result. With environments at $\env$ and once closures as cells (\S\ref{sec:model-functions}), both
facts travel with the value and appear in summaries; the programs are F1 and F2 of
Table~\ref{tab:counterexamples}, now rejected. A closure \emph{stored} in a structure holds its
captures exactly as long as the structure is live, so a record of handlers built from a list of columns
keeps that list from being edited in place while the record is used, and no longer. \emph{Solved.}

\section{Pipelines, identity and composition}
\label{sec:hard-pipelines}

The pipe is syntax: \code{xs ▷ List.filter p ▷ List.set 0 v} \emph{is}
\code{List.set (List.filter xs p) 0 v}, because pipes are rewritten before anything runs, and a
placeholder or a \code{←} bind is a written lambda, so there is no plumbing function in which to lose
precision. \code{identity x = x} is inferred with $\mathit{res} = \{\pi_1\}$; a user-written
\code{apply f x = f x} with $\mathit{res}$ equal to its callback's result translated through its supply,
here $\{\pi_2\}$ wherever the callback returns its argument; a composition \code{λx → g (f x)} likewise.
Futhark's alias analysis is lossy exactly here and its author is exploring parametricity as the fix
\citep{henriksen2026aliasing}; we analyse the body and let the class carry the result's provenance, so
identity, application and the pipe keep exact aliasing with no free theorem needed. \emph{Solved.}

The residual cost lay in the identity promises of beni's current list library
(\S\ref{sec:bg-lists}): \code{filter} that keeps every element returns its input itself, as does
\code{map} when every result is \code{===} its element. Under the rule such a function has
$\mathit{res} = \{\fresh, \pi_1\}$---the result \emph{may} be the input---so
\code{List.set (List.filter xs p) 0 v} would demand \code{xs} unique too. We propose dropping the
promises that alias (\S\ref{sec:model-demands}; \S\ref{sec:disc-options}, Change~2), after which a
pipeline never demands its source.

\section{Generic code}
\label{sec:hard-generic}

A type's \code{eq} and \code{compare} only read, by the rule of \S\ref{sec:inf-generic}, so
\code{where k.compare : k, k → Order} and \code{a.eq} need no class. A user method's summary is a class
on the constraint's method type, bound at each call site from the dispatch table's term. Inside the
generic body, a call \code{x.m a} is a call through a function-typed value whose class is the
constraint's, and the body's own summary is expressed over it. A generic body cannot write a value of
type \code{a} in place, because there is no demand of type \code{a → a}, which is the right answer.
\emph{Solved}, at the cost of one rule on \code{eq} and \code{compare}.

\section{Lists in containers}
\label{sec:hard-containers}

\paragraph{A record.} In \code{\{ m | rows = List.push m.rows x \}} the operand is $(m, .\mathit{rows})$,
and the holders of that cell are $(m, .\mathit{rows})$ itself and whatever else holds \code{m} or its
\code{rows}. A record update shares every \emph{other} field's value between the old and the new
record, but not \code{rows}, which the push's result replaces. After the update \code{m} is dead unless
used later, and a later \code{m.name} reads a path that does not meet $.\mathit{rows}$, so it is not a
use of the cell. A helper that receives the model and replaces \code{rows} is $\unused$ at
$\pi_1.\mathit{rows}$, so passing the old model to it after the push is not a use either.
\emph{Solved}: a model threaded through helpers that each touch one field is accepted, because
$\mathit{use}$ and $\live$ are per path.

\paragraph{A tuple, a constructor.} The same, through $.i$ and $C\#i$.

\paragraph{A list of lists.} In \code{List.update rows i (λr → \{ r | tags = List.push r.tags t \})} the
lambda's parameter \code{r} is an element of \code{rows}. Two slots of \code{rows} may hold the same
record (\code{[ r, r ]}), or one record may hold one list at two fields; in either case a push through
one path would be seen at the other, and the cell-set domain cannot tell slots apart
(\S\ref{sec:model-excl}). The derivation therefore needs two things: \code{List.update} must be a
consuming writer that overwrites slot $i$ with the callback's result---a destructive read, so the slot's
old holder is dead during the callback---and \code{rows} must be unique and exclusive, in the deep sense
of \S\ref{sec:model-excl}. Then \code{r} is unique inside the callback, its paths are distinct cells,
\code{r.tags}'s only holder is \code{r}, and the push is in place: accepted. If \code{rows} is unique but
not exclusive---its elements came from a \code{foreign} that declares nothing, from \code{List.repeat},
or from a literal that placed one list at two fields---the program is rejected: as a limit, ``the
elements of \code{rows} may share'', when the checker cannot see the construction, and as sharing when
it can; the fix is \code{List.copy r.tags}. With \code{map} in place of \code{update}, \code{map} supplies
a borrowed element---the old \code{rows} holds every element until \code{map} has built its result---so
the push is rejected as real sharing. This one example contains Lean's motivating \code{groupBy},
Clean's propagation rule and the two-slots-one-record hazard. \emph{Solved with a cost}: a nested edit
needs exclusivity, which the program's own construction of the container must establish, and which the
evaluation counts.

\paragraph{A dictionary.} beni's \code{Dict k v} is a balanced tree of constructors written in beni with
no \code{Js}, and it is exported opaquely. Its paths fold (\S\ref{sec:model-cells}): every value of the
tree is reached by one folded path, which importers see as the opaque step $\langle\mathit{Dict}\rangle$
and which we write $.\mathit{contents}$. Its summaries are therefore inferred, not asserted. A lookup
\code{Dict.get d k} is a non-destructive read: its result aliases a value inside \code{d}
($\mathit{res} = \{\pi_1.\mathit{contents}\}$), and \code{d}'s values stop being exclusive while that
result is live. \code{Dict.insert d k v} stores \code{v} at $.\mathit{contents}$ of its result, rebuilds a
path and shares the rest with \code{d}; it keeps $\excl$ only if \code{v} was stored exclusively.
\code{Dict.update d k f}, a path-copying walk that takes the node's value, calls \code{f} and rebuilds the
node with the result, is a destructive read at a key, which the folded path cannot name, so its
callback's owned supply is the one asserted row of the module. Huisinga's motivating example for Lean
\citep[\S2.1]{huisinga2023lean}, in which a group is looked up, pushed to and re-inserted, is
\emph{rejected} as real sharing with the holder named---the dictionary is live across the push,
because it is passed to \code{insert} afterwards---and the message offers \code{Dict.update}, which is
accepted when the values are exclusive (Appendix~\ref{app:groupby}). Lean accepts the lookup form and
copies silently at run time, because the dictionary still holds the group; the program is quadratic
\citep[\S2.4]{huisinga2023lean}. \emph{Solved with the same cost.} A set is the same.

\section{Recursion and folds}
\label{sec:hard-recursion}

\begin{Verbatim}
go xs acc = case xs of
    [] → acc
    [ x, …rest ] → go rest (List.push acc (f x))
\end{Verbatim}
\code{go}'s group is one declaration. The optimistic start assumes \code{acc} $\unused$ and the result
empty. In round 1 the tail call passes \code{List.push acc ⋯}, a demand on \code{acc}'s cell, whose only
holder in the frame is \code{acc} itself; after the jump \code{acc} is rebound, so it is dead. Hence
$\mathit{use}(\mathit{acc}) = \{\mathsf{R}, \mathsf{A}, \mathsf{W}\}$, and $\mathit{res} = \{\pi_2\}$ (the
\code{[]} arm returns \code{acc}) united with the recursive call's result. Round 2 gives the same;
validation passes. \code{List.foldl} in the core library has this shape, with its callback's class
carrying the accumulator's ownership (\S\ref{sec:inf-classes}), so
\code{List.foldl xs [] (λx acc → List.push acc x)} is accepted and runs as the in-place builds of our
prototype study did. Mutual recursion is the same fixpoint over the group. \emph{Solved.}

\section{Branches that keep the old value}
\label{sec:hard-branches}

\begin{Verbatim}
case List.get xs i of
    Just v  → v
    Nothing → List.set xs i d ▷ List.get ⋯
\end{Verbatim}
\code{v} is an element, not the array, and the \code{set} writes a slot; per-arm liveness makes
\code{xs} dead after the \code{set} on the \code{Nothing} arm, where it is the operand. Accepted. The
shape that Rust's borrow checker has long rejected and still rejects on its stable release---the third
``problem case'' of non-lexical lifetimes \citep{rfc2094}, deferred from that work to Polonius
\citep{matsakis2018alias}---does not arise, because beni has no references into a list: \code{get}
returns the element, not a pointer to the slot.

In \code{ys = if c then List.set xs 0 1 else xs}, $A(\mathit{ys})$ is \code{xs}'s cell on both arms---in
the \code{then} arm as the demand's result, in the \code{else} arm as an alias. The demand in the
\code{then} arm asks whether \code{xs} is live after it \emph{on that arm}, and it is not: the
\code{else} arm's use of \code{xs} is not on the \code{then} path. Accepted, and \code{ys} owns the cell
on both arms. Futhark's rule that the aliases of a conditional are the union of both branches
\citep[Fig.~5]{henriksen2017futhark} is kept for cells, but liveness is per path, so a value ``used''
only in the arm not taken causes no false error. \emph{Solved.}

What remains coarse: in \code{xs = if c then a else b} followed by \code{List.set xs 0 1}, both \code{a}
and \code{b} must be unique, so a later read of \code{a} on the path where \code{xs} is \code{b} is a
false error. A known imprecision; the evaluation counts how often it occurs.

\section{The foreign boundary}
\label{sec:hard-foreign}

A \code{foreign} has no body; its declaration carries a summary as it already carries an effect rung
and its \code{sync} marks. \emph{The default is pessimistic}: a parameter that says nothing is read and
escaped, because beni's existing contract lets a sibling keep a list it was given in order to read it
later, as some platform code does, and a result that says nothing is $\top$. A sibling that only reads
its argument during the call says borrowed; one that returns a fresh plain array that nothing else
holds (already the rule for arrays a sibling returns) says fresh; \code{Debug.log}, which is beni over
\code{Js}, has an asserted row saying that it reads its second argument and returns it. A platform
author who writes nothing therefore loses precision and never soundness. \code{Js.from x} marks
\code{x} escaped forever; a \code{Js.call} that passes a list is a call with no summary, which reads and
escapes it. The existing build-time checks on siblings are unchanged; the new promise is tested as the
effect rung is, by review and by a test program that pins each row (A8). The residual risk is a sibling
that lies---a bug in privileged code of the same class as a wrong effect rung. \emph{Solved, by
contract} (A2), with the trusted surface of \S\ref{sec:soundness-trusted}.

\section{Fibers and commands that capture values}
\label{sec:hard-fibers}

\code{Cmd.perform (λsend → ⋯ model.todos ⋯)} captures \code{model.todos} in a thunk that the runtime
stores and runs later. \code{Cmd.perform}'s row puts $\mathsf{E}$ on its argument's $\env$, so every
capture of the lambda is escaped at the capture, and \code{model.todos}'s cell is escaped from then on,
in this arm \emph{and in every later dispatch that receives the same cell in its model}
(\S\ref{sec:hard-tea}). Spawning a fiber (in the current scope or a named one), the thunks of a
parallel combinator and \code{Ref.set} are the same. A thunk the callee calls synchronously and does
not keep---the body of a structured-concurrency scope, the acquire step of a bracket---is called at most
once and kept by nothing; a bracket's \emph{release} is stored on the fiber's finaliser list and
escapes. These facts are asserted rows of the core \code{Task} module and inferred for its beni
wrappers.

A command whose thunk cannot suspend is a special case worth stating. Its body runs to completion in
one go; whether it can still observe a later in-place write depends only on whether the runtime runs it
before the next dispatch. Today's TEA runtime queues it to run soon after \code{update} returns, without
promising that it runs before the next message, so its captures are escaped. A platform that promises
to drain such commands before the next dispatch (\Obl{P7}) may declare their captures held only until
then, and the entry protocol of \S\ref{sec:hard-tea} uses that: a sync command that encodes and stores
the todo list it captured no longer makes the list uneditable. The checker knows which thunks cannot
suspend from the effect bits it already infers. \emph{Solved by escape}, with the cost that a list
handed to a command that may wait, or to a fiber, is shared until the end of the program, because the
analysis cannot see when the fiber reads it; the message names the capture and says that this is
\emph{possible} sharing (\S\ref{sec:errors-examples}).

\section{The TEA runtime: an entry protocol}
\label{sec:hard-tea}

A TEA platform calls the program's functions---\code{init}, \code{update}, \code{view}, subscriptions,
the bodies of commands, and, on the direct platform, derived values and listener bodies---on a model the
runtime holds between calls. None of these calls has a caller in the program to discharge \Unescaped,
\Dead\ and \Disjoint\ for the model, so the platform must say, for each one, what it passes owned, what
it passes only to read, and what it keeps of the result. An earlier version of these rules stated such a
protocol for \code{update} alone. It must cover every entry, or a \code{view} that writes a model list
in place is accepted and writes it at every render.

\paragraph{The protocol.} Table~\ref{tab:entry} gives the protocol we propose, for both of beni's
browser platforms. It is an asserted row on the platform's program constructor (\code{Tea.element} and
its kin), whose record parameter's fields carry the classes in the table, and it models the runtime's
loop---the model is \code{init}'s result, then each dispatch's result---as a recursive caller of
\code{update}. The checker evaluates that row \emph{at the one call that starts the program}, from the
published summaries of the functions in the record, as it evaluates any call: it computes the abstract
model at entry as the least fixpoint of \code{init}'s result and \code{update}'s result over it, with
numbered fresh cells for aliasing and $\mathsf{E}$ for escapes, and then checks every path that
\code{update} writes as a caller would check an argument. Nothing in this is whole-program: it reads
interfaces.

\begin{table}[t]
\centering\small
\begin{tabularx}{\linewidth}{@{}>{\raggedright\arraybackslash}p{0.22\linewidth}>{\raggedright\arraybackslash}p{0.3\linewidth}X@{}}
\toprule
Function & Receives & What the runtime keeps of the result \\
\midrule
\code{init ⊤} & nothing & the model, handed to the first \code{update}; commands, as below \\
\code{update msg model} & \code{model} owned (\Obl{P1}); \code{msg} owned unless a view event built it from the model (\Obl{P6}) & the new model, handed to the next \code{update} (\Obl{P1}); commands, as below \\
\code{view model} & \code{model} to read only & template runtime: the rendered view, including every list it shows and every handler closure, \emph{escaped}; direct platform: nothing, since \code{view} is compiled away \\
a derived value (direct) & the model, to read only & its slot; a list-valued slot is recomputed before it is read (\Obl{P4}) \\
a listener body (direct) & the current model and row item, to read only (\Obl{P3}) & the message it builds \\
\code{subscriptions model} & \code{model} to read only & the subscriptions and the functions in them, \emph{escaped} \\
a command thunk & its captures & thunks that may wait: \emph{escaped}; thunks that cannot suspend: held until the next dispatch if the platform claims \Obl{P7}, else escaped \\
\bottomrule
\end{tabularx}
\caption{The entry protocol of a TEA platform.}
\label{tab:entry}
\end{table}

\paragraph{The obligations.} The table rests on seven obligations on the platform's lowering and
runtime.
\begin{description}[leftmargin=1.2em,font=\normalfont\itshape]
\item[\Obl{P1}, hand-over.] The runtime calls \code{update} with the model \emph{owned}: after the call,
it holds no content reference to any list reachable from the old model, and the module-level model
variable is reassigned to the result.
\item[\Obl{P2}, reads before the arm are the handler's.] Anything the runtime needs from the old
model---the indices it uses to locate rows, a derived value's old input---is read \emph{before} the arm,
into locals or as scalars. A handler that read the old list and compared it after the arm would commit a
sharing error, which the lowering must not generate.
\item[\Obl{P3}, reads by listeners and items.] Row items in instance arrays are compared only by
\code{===} (A4). A listener body runs when its event fires, as program code, on the model and row item
that are current \emph{at that moment}: after every dispatch the runtime re-obtains each instance's item
from the new model for every row whose path is in the write set, so that no listener uses a reference
obtained before an in-place write. Rows outside the write set hold items whose cells were not written.
When a row's item is itself a list, its row is re-patched when its path is in the write set even though
its identity is unchanged (A4).
\item[\Obl{P4}, slots are recomputed.] A derived value's slot that holds a list is recomputed by every
message whose write set conflicts with its reads, before any DOM write reads it, and the old slot value
is never dereferenced after the arm. A derived value only reads the model, so a derived value that
writes a model list in place is an error at its definition.
\item[\Obl{P5}, the entry is checked as a call.] The abstract model at entry is the fixpoint described
above. \emph{Aliasing between paths}: an \code{init} that places one cell at two paths, \code{xs = ⋯;
\{ a = xs, b = xs \}}, publishes $\mathit{res}(.a) = \mathit{res}(.b) = \{\fresh_1\}$, so a push on
\code{model.a} fails \Disjoint\ against \code{model.b} at the program's start. \emph{Escapes across
dispatches}: a cell that \code{update}'s summary marks $\mathsf{E}$---captured by a command, a
subscription, a stored closure, a \code{Js} call or a send---is escaped at every later entry, because the
runtime may hand the same cell back in the next model.
\item[\Obl{P6}, payloads.] A message's payload list is unique at the arm only if nothing else holds it
when the arm runs. A program sender---a fiber calling \code{send (Got rows)}---gives it up: sending
consumes the message (\S\ref{sec:disc-options}, Change~9), so a sender that uses \code{rows} afterwards
is rejected at the sender. A \emph{view event} whose message is built from a model path or a row item
(\code{onClick=\{Keep model.rows\}}) hands the arm a cell the model still holds; such a payload path is
aliased with the model at the arm, and the message offers to copy it in the arm or to send an
identifier instead of the list.
\item[\Obl{P7}, sync commands are drained.] Optional: a platform may promise that a command whose thunk
cannot suspend runs to completion before the next message reaches \code{update} (A12). Without the
promise, such commands' captures are escaped like any other command's.
\end{description}

\paragraph{The two platforms.} beni's current template runtime keeps the rows it rendered and the last
value of every hole (\S\ref{sec:bg-browser}); its row says that \code{view}'s result is escaped, so every
list \code{view} shows is escaped from the first render, and an in-place edit of it in \code{update} is
an error on every message. That is the platform both applications we examined run on
(\S\ref{sec:disc-apps}); on it, the rule hosts in-place edits only of model lists that are never
rendered. A refinement---escaping only the model paths that \code{view}'s summary passes to markup---would
recover some of those, and is not designed here. The direct platform removes these holds by design: it
has no \code{view} at run time and no retained render, slots hold leaf values or recomputed lists,
instance arrays hold row items compared by identity, edit scripts are guarded by their own tags and
never by list identity, and the old model is dead after the arm. Under \Obl{P1}--\Obl{P6} the
derivation for a table application's row swap and row update, and for the in-place \code{List.update}
and \code{List.push} of the compiler's test programs for that platform, is: the model is owned at entry;
\code{model.rows} is not escaped by any arm and is not aliased by another path; the operand's only
holder is $(\mathit{model}, .\mathit{rows})$; the model is dead after the arm. The update is accepted,
and the write is \code{a[i] = v} on a plain array, derived from the general rule rather than a
model-specific one. \emph{Solved, conditional on the platform keeping the protocol}; \Obl{P2},
\Obl{P3}'s re-obtaining of items, and \Obl{P7} are not yet in the direct platform's specification.

\paragraph{\code{init} is a function.} In beni's TEA platforms today \code{init} is a value, and
applications declare it at top level. A top-level value is a holder for the program's life
(\S\ref{sec:model-check}), so every cell \code{init} makes is escaped, and the first in-place edit of a
list that started in \code{init} is rejected. That is the sound reading, and the alternative---letting
the protocol ignore the top-level hold---is unsound: a program whose \code{Reset} message returns
\code{init} would get back the array that earlier messages had written. We therefore propose that
\code{init} become a function, \code{init : ⊤ → Model × Cmd Msg}, as Elm's own \code{init} is a function
of its flags (\S\ref{sec:disc-options}, Change~12). The runtime calls it once; a \code{Reset} arm calls
it again and gets fresh cells. A program that keeps \code{init} as a value loses nothing in
correctness; its lists that start in \code{init} cannot be edited in place.

\paragraph{\code{view}, derived values and spreads.} \code{view} receives the model to read. Under the
earlier version of these rules a spread or \code{++} was a demand, so
\code{shown = [ …model.pinned, …model.recent ]} in a \code{view} wrote \code{model.pinned} in place at
every render, or, if \code{view} was declared read-only, was an error. With spreads building new arrays
(\S\ref{sec:model-demands}) it is neither: it reads two model lists and builds a third. A \code{view},
a derived value or a subscription that calls a named writer on a model list is an error at its
definition, and the message says that the runtime calls it again and again on the same model.

\section{The explicit copy}
\label{sec:hard-copy}

\code{List.copy : List a [borrowed] → List a [= fresh]} is a shallow copy, \code{a.slice()}, or for a view
\code{b.slice(o)}. The elements are the same values; records are immutable and need no copying; a
list nested inside an element is still shared, and editing it in place is still an error whose message
names the inner holder, with the fix \code{List.map xs List.copy} or a copy at the inner edit. A copy has a
new identity, so a list keyed by reference in the view sees new rows, which is right for a list whose
contents the program is about to change behind the original's back. Copying a list that is already
unique is legal and wasteful; the evaluation counts such copies to decide whether a warning is
wanted---a warning, never an error. \emph{Solved.}

\section{Views}
\label{sec:hard-views}

\code{[ x, …rest ]} binds \code{rest} to a view $(b, o)$ of \code{xs}'s base array; \code{List.drop} and
\code{List.tail} return views too. A view is a holder of the base, and a write through a view writes the
base: \code{List.set rest i v} is \code{b[o + i] = v}. The demand on \code{rest} therefore asks that
every holder of the base---\code{xs} included---be dead. In the common walk (\code{go rest ⋯} with
\code{xs} not used again) it is; where \code{xs} is used afterwards, the rejection is correct, because
the elements \emph{are} shared. A view is always a suffix of its base, so \code{push} on a view is a
push on the base. A view header also caches a plain copy of its elements and stores its length; a write
through a view therefore never mutates the header it was given, but returns a new header with no cached
copy and the right length, and the old header, being the consumed operand, is dead (A1). A reader that
took the cached copy earlier holds a separate array, a snapshot of the old version, which is what value
semantics gives it. The scalar-view rewrite removes most views from loops anyway (\S\ref{sec:bg-lists}).
\emph{Solved.}

\section{Mutable cells}
\label{sec:hard-ref}

\code{Ref.get r} returns an alias of whatever the cell holds, and \code{Ref.set r xs} stores \code{xs};
the analysis cannot see who else has called \code{Ref.get}. A list stored in a \code{Ref} is therefore
escaped, every \code{Ref.get} result is $\top$-held, and \code{Ref.update r (λxs → List.push xs x)} is
rejected as a limit, with a message saying that the checker cannot see the cell's other readers. A
linear-cell discipline---the cell is the only holder, and \code{update}'s callback consumes its
argument---would need a proof that no \code{Ref.get} result is live across the update, a
whole-program fact about an unbounded set of readers. The fix is a copy in the callback, or a
persistent structure in the cell. A queue, if one is added to the core library, is the same.
\emph{Open, as a precision limit}; the evaluation measures how often it fires.

\section{Messages as values}
\label{sec:hard-messages}

A message stored in the model, logged whole, or mapped by \code{Html.map} with a function that is not a
constructor is a value; its payload lists are held wherever the message is. A payload passed as
handler parameters---the common case---is covered by \Obl{P6}. \emph{Solved.}

\section{Versions: undo histories, diffs, snapshots}
\label{sec:hard-undo}

In \code{\{ model | history = [ model.todos, …model.history ], todos = List.push model.todos t \}} the
record update evaluates its fields in written order, so \code{history} holds the old array when the
push runs---a live holder, used when the record is built after the push. The program is
\emph{rejected}, correctly: the undo entry would change with the push. The message says so and offers
\code{List.copy model.todos} in the history, $O(n)$ per edit, the price of a version---or the
non-writing \code{model.todos ++ [ t ]} in place of the push, which builds the new version and keeps the
old. A persistent vector written in beni as a library module---the current trie, with no in-place
writes---would give $O(\log n)$ versions to programs that want them; it would be a library, not the
representation of \code{List}. \emph{Solved}, with the cost stated.

\section{Debugging aids}
\label{sec:hard-debug}

\code{Debug.log tag x} reads \code{x} at the call and returns it; \code{Debug.toString} reads at the call.
Neither is a holder afterwards; both are asserted rows (\S\ref{sec:soundness-trusted}). A development
crash screen must print only the exception and the stack (A9); one that printed the model would read a
half-updated record, and the two semantics would differ \emph{after} the program has stopped---a
difference in a debugging aid, not in the program, but one to state. \emph{Solved}, with A9 recorded.

\section{Inlining}
\label{sec:hard-optimiser}

The release optimiser inlines a private function called from one site into that site, and folds a
single-use binding into its use (\S\ref{sec:bg-release}); it does not specialise functions per call
site. Inlining substitutes a body at a call, which preserves evaluation order and holders: the
inlined body's demands are the same demands at the same points, with the parameters bound to the
argument temporaries. The verdict was computed on the source and is not recomputed, so the accepted set
is the source's, never the optimiser's; a program the inliner could have proved safe is still rejected.
The fold is the subject of A5. \emph{Solved by construction.}
